\documentclass[11pt,a4paper]{article}
\usepackage[margin=0.55in, top=0.5in, bottom=0.5in]{geometry}
\usepackage[utf8]{inputenc}
\usepackage[T1]{fontenc}
\usepackage{charter}
\usepackage{fontawesome5}
\usepackage{xcolor}
\usepackage{titlesec}
\usepackage{enumitem}
\usepackage{tikz}
\usepackage{hyperref}
\usepackage{tcolorbox}
\tcbuselibrary{skins, breakable}
\usepackage{microtype}

\definecolor{headerbg}{HTML}{FFFFFF}
\definecolor{accent}{HTML}{1E3A5F}     % Professional Dark Navy
\definecolor{accentblue}{HTML}{2B6CB0} % Slate Blue
\definecolor{lightgray}{HTML}{F3F4F6}
\definecolor{darktext}{HTML}{1F2937}
\definecolor{skillbg}{HTML}{374151}
\definecolor{bulletline}{HTML}{9CA3AF}
\definecolor{subtletext}{HTML}{4B5563}

\hypersetup{colorlinks=true, linkcolor=accentblue, urlcolor=accentblue}

\titleformat{\section}{\large\bfseries\color{accent}}{}{0em}{}[\vspace{1pt}\titlerule]
\titlespacing{\section}{0pt}{11pt}{7pt}

\newtcbox{\skillbadge}{
    on line, box align=base,
    colback=skillbg, colframe=skillbg,
    colupper=white,
    fontupper=\small\bfseries,
    arc=4pt, boxrule=0pt,
    left=6pt, right=6pt, top=3pt, bottom=3pt,
    boxsep=0pt
}

% Caja de entrada de experiencia con borde izquierdo elegante y continuo
\newtcolorbox{timelineentry}{
    blanker,
    breakable,
    borderline west={1.5pt}{0pt}{bulletline},
    left=1.2em,
    top=0pt,
    bottom=0pt,
    before skip=0.8ex,
    after skip=2.6ex
}

\newcommand{\cventry}[5]{%
  \noindent
  \begin{minipage}{\textwidth}
    \vspace{0pt}%
    \tikz[baseline=-0.7ex]{\node[draw=accent, fill=white, circle, inner sep=2.5pt, line width=1.5pt] {};}%
    \hspace{0.5em}%
    {\large\bfseries\color{accent} #1}%
    \hfill{\small\bfseries\color{accentblue} #4}\par
    \vspace{0.2ex}%
    \hspace{1.5em}{\normalsize\bfseries\color{darktext} #2}%
    \hfill{\small\color{subtletext}\textit{#3}}\par
  \end{minipage}\par
  \vspace{0.4ex}%
  \begin{timelineentry}
    \small\color{darktext}
    \begin{itemize}[label={\color{accentblue}\raisebox{0.2ex}{\tiny\faCircle}}, leftmargin=1.2em, itemsep=0.4ex, parsep=0pt, topsep=0pt, partopsep=0pt]
      #5
    \end{itemize}
  \end{timelineentry}
}

\begin{document}
\thispagestyle{empty}
\color{darktext}

% Encabezado
\begin{center}
    {\Huge \bfseries\color{accent} SANTIAGO ALEJANDRO OCHOA QUESADA} \\[6pt]
    {\small\color{darktext}
    \faEnvelope\ \href{mailto:santy.alejo4@gmail.com}{santy.alejo4@gmail.com}
    \quad\textcolor{bulletline}{\textbar}\quad
    \faPhone\ +57 320 351 7153
    \quad\textcolor{bulletline}{\textbar}\quad
    \faMapMarker*\ Bogotá, Colombia \\[4pt]
    \faLinkedin\ \href{https://linkedin.com/in/santiago-alejandro-ochoa-quesada-a07a18288/}{LinkedIn}
    \quad\textcolor{bulletline}{\textbar}\quad
    \faGithub\ \href{https://github.com/Sochoaq}{GitHub}
    }
\end{center}
\vspace{4pt}

\section{Perfil Profesional}
\textbf{Ingeniero Electrónico} de la Universidad Nacional de Colombia (promedio 4.6/5.0) y actual \textbf{Ingeniero de Proyectos} en AM4. Especializado en ingeniería de software, integración de datos (ETL), desarrollo backend en Python y gestión de bases de datos relacionales SQL. Experiencia en el ciclo de vida completo de aplicaciones, diseño de APIs REST, contenedores Docker, pruebas automatizadas (QA con Pytest) y modelos de Machine Learning. Con sólida formación en redes, ciberseguridad e infraestructura bajo estándares de alta disponibilidad. Residencia en Bogotá con disponibilidad para modalidad híbrida o remota a nivel nacional. Nivel de inglés profesional B2.

\section{Experiencia Profesional}

\cventry{Ingeniero de Proyectos (Software y Automatización)}{AM4}{Bogotá}{04/2026 - Actualidad}
{
    \item Diseño, desarrollo e implementación de flujos automatizados de integración de datos (ETL) en Python y SQL para la captura, limpieza y transformación de archivos y telemetrías operativas.
    \item Construcción de microservicios y APIs REST para conectar fuentes de datos heterogéneas, bases de datos relacionales y módulos analíticos de procesamiento continuo.
    \item Implementación de buenas prácticas de ingeniería de software, arquitectura de datos y modelado entidad-relación, asegurando consistencia y trazabilidad de los flujos.
    \item Ejecución de pruebas unitarias y validación funcional (QA con Pytest), alcanzando un 98\% de confiabilidad operativa en implementaciones críticas para clientes industriales (Refoenergy, Brinsa, Geopark).
    \item Contenedorización de aplicaciones y servicios mediante Docker, facilitando el despliegue estandarizado, la observabilidad y la estabilidad operativa.
    \item Documentación técnica integral de sistemas: diagramas de casos de uso, especificación de APIs, modelos de datos y reportes de desempeño operativo.
    \item Participación activa en la resolución de anomalías, soporte técnico y coordinación de entregables técnicos bajo metodologías colaborativas e iterativas.
}

\cventry{Auxiliar Estudiantil}{Universidad Nacional de Colombia}{Bogotá}{07/2024 - 11/2025}
{
    \item Apoyo en laboratorios de ingeniería mediante simulación computacional, programación de algoritmos y análisis y validación de datos experimentales.
    \item Preparación de registros, verificación de instrumentación técnica y protocolos de medición para asegurar la consistencia y calidad de la información.
    \item Acompañamiento a estudiantes en diagnóstico de fallas, depuración de software y resolución estructurada de problemas analíticos.
    \item Elaboración de guías técnicas, documentación y reportes de prácticas estandarizadas para asegurar la trazabilidad del conocimiento.
}

\section{Educación}

\cventry{Ingeniería Electrónica}{Universidad Nacional de Colombia}{Bogotá}{02/2021 - 11/2025}
{
    \item \textbf{Promedio académico: 4.6 / 5.0} (Graduado con Mérito Académico).
    \item Formación sólida en sistemas computacionales, análisis numérico, procesamiento de datos, comunicaciones, automatización y control.
    \item Desarrollo de proyectos de ingeniería orientados a algoritmos aplicados, integración de software y hardware, y modelado matemático.
}

\newpage

\section{Proyectos Relevantes}

\cventry{Pipelines de Integración de Datos \& Machine Learning}{Python, SQL \& Analítica}{Bogotá}{2026}
{
    \item Desarrollo de un motor modular para la ingestión, transformación (ETL) y almacenamiento estructurado de datos masivos y archivos multi-formato, optimizando tiempos de consulta en un 35\%.
    \item Integración de modelos predictivos de Machine Learning y procesamiento estadístico para la detección temprana de inconsistencias y anomalías operativas.
    \item Construcción de endpoints REST para consulta segura de datos, control de versiones y generación automática de reportes analíticos con trazabilidad integral.
}

\cventry{DPS Analyzer \& Metrology Platform (Sistemas y Resiliencia)}{AM4}{Bogotá}{2026}
{
    \item Suite de telemetría y diagnóstico para infraestructura crítica: captura de señales a 50 kHz, procesamiento FFT y evaluación de índices operativos bajo estándares internacionales.
    \item Contenedorización completa con Docker y automatización del pipeline de construcción y pruebas CI/CD mediante GitHub Actions, garantizando despliegues repetibles y confiables.
    \item Implementación de esquemas de persistencia en bases de datos relacionales SQL, monitoreo de métricas operativas y paneles de visualización interactivos.
}

\section{Habilidades}
\begin{tcolorbox}[colback=lightgray, colframe=lightgray, boxrule=0pt, arc=6pt, left=10pt, right=10pt, top=10pt, bottom=10pt]
\noindent
\linespread{1.8}\selectfont
\skillbadge{Python (Intermedio/Avanzado)}
\skillbadge{SQL \& Modelado Relacional}
\skillbadge{Pipelines de Datos (ETL)}
\skillbadge{FastAPI / REST APIs}
\skillbadge{Docker \& Contenedores}
\skillbadge{AWS (EC2, S3, IAM, Lambda)}
\skillbadge{Pruebas QA \& Pytest}
\skillbadge{Machine Learning Aplicado}
\skillbadge{Git \& GitHub Actions CI/CD}
\skillbadge{Linux SysAdmin}
\skillbadge{Redes \& Cisco CCNA}
\skillbadge{Ciberseguridad \& Ethical Hacking}
\skillbadge{Metodologías Ágiles (Scrum)}
\skillbadge{Resolución de Problemas}
\skillbadge{Trabajo Colaborativo}
\end{tcolorbox}

\section{Idiomas}
\begin{itemize}[label={}, leftmargin=0.3cm, itemsep=6pt]
    \item \textbf{Español} \hfill Nativo \quad \textcolor{accent}{\faCircle \faCircle \faCircle \faCircle \faCircle}
    \item \textbf{Inglés} \hfill B2 (Profesional de Trabajo / Intermedio) \quad \textcolor{accent}{\faCircle \faCircle \faCircle \faCircle} \textcolor{bulletline}{\faCircle}
    \item \textbf{Francés} \hfill A2 \quad \textcolor{accent}{\faCircle \faCircle} \textcolor{bulletline}{\faCircle \faCircle \faCircle}
\end{itemize}

\section{Certificaciones}
\begin{itemize}[label={\color{accentblue}\faAward}, leftmargin=0.6cm, itemsep=6pt]
    \item \textbf{Ethical Hacking \& Cybersecurity} --- Cisco Networking Academy \hfill \textit{04/2026}
    \item \textbf{CCNA: Introduction to Networks} --- Cisco Networking Academy \hfill \textit{07/2025}
    \item \textbf{Análisis Computacional de Datos Ambientales} --- Universidad Nacional de Colombia \hfill \textit{09/2025}
    \item \textbf{Introducción a la Fotónica Integrada y las Comunicaciones Cuánticas} --- UNAL \hfill \textit{07/2024}
    \item \textbf{Python Intermedio} --- SoloLearn \hfill \textit{03/2021}
\end{itemize}

\end{document}
